\begin{quadro}[t]
    \centering
    \caption{Exemplo de cenários de teste estruturados em formato BDD (sintaxe Gherkin)}
    \label{qua:exemplo_bdd}
    
    \fcolorbox{black}{backcolour}{% 
        \begin{minipage}{0.9\textwidth}\small%
            \vspace{3pt}%
            \noindent\normalsize{\textcolor{blue}{Cenário}: Saque com saldo suficiente}\vspace{3pt}\\
            \hspace*{1.25cm}\textcolor{red}{DADO} que a conta possui saldo disponível\\
            \hspace*{1.25cm}\textcolor{red}{E} que o cartão inserido é válido\\ 
            \hspace*{1.25cm}\textcolor{red}{E} que o terminal possui cédulas suficientes\\
            \hspace*{1.25cm}\textcolor{red}{QUANDO} o cliente solicita o resgate do dinheiro\\
            \hspace*{1.25cm}\textcolor{red}{ENTÃO} a conta deve ser debitada no valor correspondente\\
            \hspace*{1.25cm}\textcolor{red}{E} o dinheiro deve ser dispensado\\
            \hspace*{1.25cm}\textcolor{red}{E} o cartão deve ser ejetado\\
            
            \noindent\normalsize{\textcolor{blue}{Cenário}: Saque com saldo insuficiente}\vspace{3pt}\\
            \hspace*{1.25cm}\textcolor{red}{DADO} que a conta não possui saldo disponível\\
            \hspace*{1.25cm}\textcolor{red}{E} que o cartão inserido é válido\\ 
            \hspace*{1.25cm}\textcolor{red}{E} que o terminal possui cédulas suficientes\\
            \hspace*{1.25cm}\textcolor{red}{QUANDO} o cliente solicita o resgate do dinheiro\\
            \hspace*{1.25cm}\textcolor{red}{ENTÃO} a transação deve ser recusada\\
            \hspace*{1.25cm}\textcolor{red}{E} o saldo da conta deve permanecer inalterado\\
            \hspace*{1.25cm}\textcolor{red}{E} uma mensagem de saldo insuficiente deve ser exibida\\
            \hspace*{1.25cm}\textcolor{red}{E} o cartão deve ser ejetado
        \end{minipage}
    }

    \vspace{5pt}
    \hspace*{0.75cm}\raggedright{\small \textbf{Fonte:} Elaborado pelo autor.}
\end{quadro}